% ============================================================
%  08_multi_runs.tex — Section 8 : Validation multi-runs
% ============================================================
\section{Validation multi-runs et classification de stabilité}

\dateblock{avril 2026}

\subsection{Motivation : non-déterminisme de xTB}

L'optimisation xTB repose sur une descente de gradient. Combinée à la perturbation aléatoire
$\pm 0.1$~\AA{} sur les coordonnées $z$ introduite plus haut, chaque exécution part d'un point
initial légèrement différent et peut converger vers un minimum local distinct. Pour les solutions dont l'angle de non-planarité est proche du seuil 10°, le
résultat d'un \textbf{run unique} n'est pas fiable : la même solution peut être marquée plane
dans 60\% des runs et non plane dans les 40\% restants.

Pour distinguer les vraies solutions plates des cas marginaux, le pipeline a été étendu pour
exécuter chaque solution \textbf{N fois} et caractériser sa \textbf{stabilité} via la
distribution des angles observés.

\subsection{Drapeau \texttt{--n-runs N}}

Le drapeau \texttt{--n-runs N} (par défaut 1) contrôle le nombre d'exécutions xTB par
solution. Chaque run utilise une nouvelle perturbation aléatoire et est sauvegardé dans un
sous-dossier dédié.

Structure de sortie pour une solution donnée :

\begin{lstlisting}[style=benzai]
output/h4/default/0-5-6-11/solutions/sol_1_5_7_7_5/
  source.xyz              # Coordonnees initiales (avant xTB)
  run_01_opt.xyz          # Resultat optimise du run 1
  run_02_opt.xyz
  ...
  run_10_opt.xyz
\end{lstlisting}

\subsection{Statistiques agrégées par solution}

Le script \texttt{aggregate\_runs.py} (appelé automatiquement par \texttt{batch\_main.py}
quand \texttt{N>1}) lit les N résultats et calcule pour chaque solution :

\begin{itemize}
  \item \texttt{n} : nombre de runs réussis (un run en échec ou en timeout n'est pas compté)
  \item \texttt{planar\_count} / \texttt{non\_planar\_count} : compteurs binaires (seuil 10°)
  \item \texttt{planar\_pct} : pourcentage de runs déclarés plans
  \item \texttt{angle\_mean}, \texttt{angle\_std} : moyenne et écart-type des angles ACP
  \item \texttt{angle\_min}, \texttt{angle\_max} : amplitudes observées
  \item \texttt{angles[]} : tableau brut des N valeurs (utile pour rejouer une analyse sans
    relancer xTB)
  \item \texttt{classification} : étiquette qualitative (cf. ci-dessous)
\end{itemize}

\subsection{Classification de stabilité}

Six classes qualitatives caractérisent la stabilité d'une solution face au non-déterminisme.
Les seuils ont été choisis pour distinguer les cas \emph{vraiment} plats des cas
\emph{borderline} :

\begin{center}
\begin{tabular}{lll}
  \toprule
  Classe & Critère & Interprétation \\
  \midrule
  \texttt{always\_planar}      & 100\% plans ET $\mu < 5°$           & vraiment plate \\
  \texttt{mostly\_planar}      & $\geq 70$\% plans ET $\sigma < 3°$  & probablement plate \\
  \texttt{unstable}            & 30--70\% plans OU $\sigma > 5°$     & cas limite (bistable) \\
  \texttt{mostly\_non\_planar} & $< 30$\% plans                      & probablement non plate \\
  \texttt{always\_non\_planar} & 0\% plans                           & vraiment non plate \\
  \texttt{ambiguous}           & autres cas                          & non concluant \\
  \bottomrule
\end{tabular}
\end{center}

\begin{remarque}
  La classe \texttt{unstable} est la plus intéressante scientifiquement : ces solutions sont
  géométriquement au seuil de la planarité. Soit la structure est réellement bistable (deux
  minima locaux distincts, l'un plat et l'autre courbé), soit la dispersion est trop grande
  et le seuil 10° n'est pas pertinent pour ce cas. Une analyse approfondie demanderait
  $N \geq 30$ runs pour confirmer la classification.
\end{remarque}

\subsection{Justification statistique du choix de N}

Pour détecter avec 95\% de confiance qu'une solution bascule au moins une fois sur 10\% des
runs, il faut $N \geq \lceil \log(0.05) / \log(0.90) \rceil = 28$. Pour des bascules à 5\%,
$N \geq 59$. En pratique :

\begin{center}
\begin{tabular}{lll}
  \toprule
  Phase & N recommandé & Objectif \\
  \midrule
  Screening rapide  & 10      & Identifier les candidats instables \\
  Analyse standard  & 20--30  & Classification fiable \\
  Analyse fine      & 50--100 & Caractérisation des cas limites \\
  \bottomrule
\end{tabular}
\end{center}

La stratégie adoptée est à \textbf{deux vitesses} : N=10 sur l'ensemble du dataset pour
identifier les cas instables, puis N=30+ ciblé sur ceux-ci si nécessaire.
